% Options for packages loaded elsewhere
\PassOptionsToPackage{unicode}{hyperref}
\PassOptionsToPackage{hyphens}{url}
\documentclass[
]{article}
\usepackage{xcolor}
\usepackage{amsmath,amssymb}
\setcounter{secnumdepth}{-\maxdimen} % remove section numbering
\usepackage{iftex}
\ifPDFTeX
  \usepackage[T1]{fontenc}
  \usepackage[utf8]{inputenc}
  \usepackage{textcomp} % provide euro and other symbols
\else % if luatex or xetex
  \usepackage{unicode-math} % this also loads fontspec
  \defaultfontfeatures{Scale=MatchLowercase}
  \defaultfontfeatures[\rmfamily]{Ligatures=TeX,Scale=1}
\fi
\usepackage{lmodern}
\ifPDFTeX\else
  % xetex/luatex font selection
\fi
% Use upquote if available, for straight quotes in verbatim environments
\IfFileExists{upquote.sty}{\usepackage{upquote}}{}
\IfFileExists{microtype.sty}{% use microtype if available
  \usepackage[]{microtype}
  \UseMicrotypeSet[protrusion]{basicmath} % disable protrusion for tt fonts
}{}
\makeatletter
\@ifundefined{KOMAClassName}{% if non-KOMA class
  \IfFileExists{parskip.sty}{%
    \usepackage{parskip}
  }{% else
    \setlength{\parindent}{0pt}
    \setlength{\parskip}{6pt plus 2pt minus 1pt}}
}{% if KOMA class
  \KOMAoptions{parskip=half}}
\makeatother
\setlength{\emergencystretch}{3em} % prevent overfull lines
\providecommand{\tightlist}{%
  \setlength{\itemsep}{0pt}\setlength{\parskip}{0pt}}
\usepackage{bookmark}
\IfFileExists{xurl.sty}{\usepackage{xurl}}{} % add URL line breaks if available
\urlstyle{same}
\hypersetup{
  pdftitle={AI Agents Development},
  hidelinks,
  pdfcreator={LaTeX via pandoc}}

\title{\texttt{AI} Agents Development}
\author{}
\date{}

\begin{document}
\maketitle

\subsection{Description}\label{description}

This two day course covers the design and implementation of autonomous
\texttt{AI} agents powered by large language models. Students will learn
to build agents that can reason, plan, use tools, and collaborate in
multi-agent systems using modern frameworks such as \texttt{LangGraph},
\texttt{CrewAI}, and \texttt{AutoGen}. The course addresses agent
memory, state management, evaluation strategies, guardrails, and
production deployment.

\subsection{Duration}\label{duration}

16 hours / 2 days

\subsection{Intended Audience}\label{intended-audience}

\begin{itemize}
\tightlist
\item
  Software developers building \texttt{AI}-powered applications
\item
  Data engineers integrating \texttt{AI} agents into pipelines
\item
  Software architects designing agent-based systems
\end{itemize}

\subsection{Prerequisites}\label{prerequisites}

\begin{itemize}
\tightlist
\item
  Proficiency in \texttt{Python} programming
\item
  Basic understanding of \texttt{LLMs} and their APIs
\item
  Familiarity with prompt engineering concepts
\end{itemize}

\subsection{Required Knowledge}\label{required-knowledge}

\begin{itemize}
\tightlist
\item
  \texttt{Python} Programming (or equivalent experience)
\end{itemize}

\subsection{Objectives}\label{objectives}

\begin{itemize}
\tightlist
\item
  Understand the architecture and \texttt{design\ patterns} of
  \texttt{AI} agents.
\item
  Build autonomous agents capable of reasoning, planning, and tool use.
\item
  Implement multi-agent systems with coordination and communication.
\item
  Apply guardrails and evaluation techniques to ensure agent
  reliability.
\item
  Deploy agents to production environments.
\end{itemize}

\subsection{Outline}\label{outline}

\begin{itemize}
\tightlist
\item
  Introduction to \texttt{AI} Agents

  \begin{itemize}
  \tightlist
  \item
    What are \texttt{AI} agents and why they matter
  \item
    Agent architectures: \texttt{ReAct}, plan-and-execute, reflexion
  \item
    Single agent vs.~multi-agent paradigms
  \item
    Overview of the agent ecosystem
  \end{itemize}
\item
  Tool Use and Function Calling

  \begin{itemize}
  \tightlist
  \item
    Defining and registering tools
  \item
    Function calling with \texttt{OpenAI}, Anthropic, and open-source
    models
  \item
    Tool selection and error handling
  \item
    Building custom tool integrations
  \end{itemize}
\item
  Agent Frameworks

  \begin{itemize}
  \tightlist
  \item
    Building agents with \texttt{LangGraph}
  \item
    Orchestrating teams with \texttt{CrewAI}
  \item
    Multi-agent conversations with \texttt{AutoGen}
  \item
    Comparing frameworks and choosing the right one
  \end{itemize}
\item
  Memory and State Management

  \begin{itemize}
  \tightlist
  \item
    Short-term vs.~long-term memory
  \item
    Conversation history and context \texttt{windows}
  \item
    Persistent state with vector stores and databases
  \item
    Stateful workflows and checkpointing
  \end{itemize}
\item
  Multi-Agent Systems

  \begin{itemize}
  \tightlist
  \item
    Agent communication protocols
  \item
    Task decomposition and delegation
  \item
    Supervisor and hierarchical patterns
  \item
    Collaborative problem-solving
  \end{itemize}
\item
  Evaluation and Testing

  \begin{itemize}
  \tightlist
  \item
    Evaluating agent performance
  \item
    Benchmarking and metrics
  \item
    Tracing and observability
  \item
    Regression testing for agents
  \end{itemize}
\item
  Guardrails and Safety

  \begin{itemize}
  \tightlist
  \item
    Input and output validation
  \item
    Content filtering and moderation
  \item
    Preventing prompt injection in agents
  \item
    Human-in-the-loop approval workflows
  \end{itemize}
\item
  Production Deployment

  \begin{itemize}
  \tightlist
  \item
    Scaling agent workloads
  \item
    \texttt{Async} and streaming execution
  \item
    Monitoring and logging in production
  \item
    Cost management and optimization
  \end{itemize}
\end{itemize}

\subsection{Copyright}\label{copyright}

Mark Veltzer
\href{mailto:mark.veltzer@gmail.com}{\nolinkurl{mark.veltzer@gmail.com}},
© 2026

\end{document}
